% Leaf bag B, binary values; every assigned factor belongs to exactly one bag.
\begin{tikzpicture}[figure picture,x=1mm,y=1mm,
  bag/.style={draw=FigureStroke,line width=.8pt,rounded corners=1.8pt,
    minimum width=51mm,minimum height=16mm,align=center,inner sep=2mm},
  every node/.append style={fill=none,inner sep=1pt}]
 \path[use as bounding box] (0,0) rectangle (169,84);
 \node[bag,fill=FigureBlueFill] (b) at (33,68) {$B=\{1,\textcolor{FigureGreen}{2,3}\}$\\局部因子组合 $\Phi_B(z_1,z_2,z_3)$};
 \node[bag,fill=FigurePurpleFill] (c) at (136,68) {$C=\{\textcolor{FigureGreen}{2,3},4\}$\\局部因子组合 $\Phi_C(z_2,z_3,z_4)$};
 \draw[FigureGreen,line width=.85pt,-{Stealth[length=1.8mm]}] (b.east)--(c.west);
 \node[above] at (84.5,69) {$S_{B,C}=\{2,3\}$};
 \node[below] at (84.5,66) {$m_{B\to C}(z_2,z_3)$};
 \node at (23,47) {保留的接口值};
 \node at (103,47) {汇总 $z_1\in\{0,1\}$ 后的消息值};
 \draw[FigureGrid,line width=.4pt] (4,42)--(165,42);
 \foreach \y/\a/\d in {35/0/0,25/0/1,15/1/0,5/1/1}{
  \node at (23,\y) {$(z_2,z_3)=(\a,\d)$};
  \node at (103,\y) {$\Phi_B(0,\a,\d)\oplus\Phi_B(1,\a,\d)$};
 }
\end{tikzpicture}
